\documentclass{article}
\usepackage[T1]{fontenc}
\usepackage{lmodern}
\usepackage{graphicx}
\usepackage{amsmath,amssymb}
\usepackage[a4paper,margin=2cm]{geometry}
\usepackage[absolute,overlay]{textpos}
\setlength{\TPHorizModule}{1mm}
\setlength{\TPVertModule}{1mm}
\usepackage[indonesian]{babel}
\usepackage{hyperref}
\setlength{\emergencystretch}{2em}
% unit: Halaman 26

\begin{document}

% === Halaman 26 (python/native) ===

























=

















































2A

F9

4F

F

3

43

BD

B

7

26


02

01

01

03

03

02

01

01

01

03

02

01

01

01

03

02

(02 • 26)  (03 • 7B)  (01 • BD)  (01 • 43) = 3F

(01 • 26)  (02 • 7B)  (03 • BD)  (01 • 43) = 4F

(01 • 26)  (01 • 7B)  (02 • BD)  (03 • 43) = F9

(03 • 26)  (01 • 7B)  (01 • BD)  (02 • 43) = 2A

Contoh:

Semua perkalian dan pernjumlahan dilakukan dalam GF(28)

Penjelasan tentang GF(2m) dapat dilihat pada lampiran

\end{document}
